%%%%% Instructions used in both pre-class and during-class portion %%%%%%

      \begin{center}{\bf Creating Functions  and Functions of Best Fit}\end{center}
  \vs{.1}
  {\large \em Using Desmos/Amplify for this Application}
  \itemC{
  \item Go to \href{https://student.amplify.com/}{\underline{student.amplify.com}} or \href{https://student.desmos.com/}{\underline{student.desmos.com}}.  
  \item You will need to  use code \desmosClassCode.
  \item You might need to register for an Amplify  account (it's free).  If so, please use your UVA email as this helps us keep track of what work your group does.
  }~\\
  {\large \em Creating Functions of Best Fit in Desmos/Amplify} \hs{.1} Functions of best fit are sometimes referred to as trend lines, regression functions or least squares lines of best fit.  The following steps will explain how to use Desmos to create functions of best fit and to find the R$^2$-value for a given set of data.
\enum{
\item {\em (This step has already been done for you.)} Go to desmos.com/calculator and create a table of data. Cut-paste from a spreadsheet is probably the most straightforward way to generate this table.  Other options are here: \red{\underline{\url{https://help.desmos.com/hc/en-us/articles/4405489674381-Tables}}} \vs{-.15}
\item {\em (This step has already been done for you.)} Find the default header for a table generated by cut-paste. It likely will be $(x_1, y_1)$. Change it to whatever you like. For example, $(t_1,c_1)$.
\item In another cell, create the linear function of best fit using variables from the header. For example
$$c_1 \sim mt_1+b$$ will create a linear best-fit function. Values for $R^2$ and   parameters $m$ and $b$ will be generated.\\
}
{\bf The following steps are not needed for pre-class activity but will be used during class.}
\enum{ \setItemnumber{4}
\item A quadratic function of best fit can be created by $\ds c_1 \sim at_1^2+bt_1+c$.  Values for $R^2$ and parameters $a$, $b$, and $c$ will be generated.
\item An exponential function of best fit can be created by  $\ds c_1 \sim a e^{b\, t_1} +c$.  Values for $R^2$ and parameters $a$, $b$, and $c$ will be generated.
\item A logarithmic function of best fit can be created by $\ds c_1 \sim a \ln(b\, t_1 +c)$. Values for $R^2$ and all parameters $a$, $b$, and $c$ will be generated.
}

    {\large \em Creating and Graphing Functions in Desmos/Amplify}
   \itemI{
   \item  Create expression cell by clicking on {\Large$+$} in the upper left corner.  Choose ``{\em $f(x)$ expression}".
   \item Type the function definition in any empty cell in left column of Desmos. For example,\\ $f(t)=3t+6$. Graphs will automatically be created while defining the function. 
   \item Function names in Desmos will utilize a single letter for the dependent variable For example, $f(t)$, $g(t)$, $h(t)$ are all acceptable function names. $Lin(t)$ or $Exp(t)$ will not be accepted by Desmos. 
   \item Subscript notation {\em can} be used in naming conventions. For example $f_{linear}(t)$ and $f_{exponential}(t)$  name two different functions.
   \item To evaluate functions at a specific value, say $t=10$,  use an empty expression cell to type in $f(1)$. The result is automatically displayed.
   }
{\large \em Data sets used in Desmos} \hs{.1}  Feel free to use these data sets in whatever tool is helpful.
\itemC{
  \item 1965 to 2010 data: \\ \red{\underline{\url{https://docs.google.com/spreadsheets/d/1KvLYR1DRHiRrkb1dh-qrSMRwzKPKuH1SkMDghl2fTcc/}}}
  \vs{-.1}
  \item 1960 to 2025 data:\\ \red{\underline{ \url{https://docs.google.com/spreadsheets/d/1dtNawrPd3W744p229R7WS83KKHVkyEokcsyCeyPDF4w}}}
  }